\documentclass{article}
\usepackage{amsmath}
\usepackage{geometry}
\usepackage{parskip}
\setlength{\parskip}{1em}
\geometry{a4paper, margin=0.8in}

\title{Week 2 Reflection}
\author{Daniyal Shahid Iqbal}

\begin{document}

\maketitle


Artificial Neural Networks (ANNs) are computational models inspired by the biological brain. The fundamental building block of Deep Learning (DL) is the \textbf{Perceptron}. The perceptron computes a linear combination of inputs and weights, which is then passed through a non-linear activation function.

Depending on the task, the output can be a probability (for classification) or a continuous number (for regression).

Since our main focus is approximating complex functions (decision boundaries), we utilize non-linear activation functions. Common functions include:
\begin{itemize}
    \item \textbf{Sigmoid:} Squashes output between 0 and 1; used for probabilities.
    \item \textbf{Tanh:} Squashes the neuron’s output between -1 and 1; strictly increasing.
    \item \textbf{ReLU (Rectified Linear Unit):} Acts as a linear function for non-negative inputs and outputs zero for negative values.
\end{itemize}


These perceptrons are stacked to form deep networks consisting of an input layer, hidden layers (where feature extraction occurs), and an output layer.

\textbf{Technical Note:} When discussing the depth of a model, we include the output layer. For example, a ``2-layer NN'' consists of one hidden layer and the final output layer.


To train these networks, we use a loss function specific to the task:
\begin{itemize}
    \item \textbf{Binary Classification:} Typically uses \textbf{Binary Cross Entropy} (the negative log-likelihood of a Bernoulli distribution), which penalizes wrong predictions.
    \item \textbf{Multinomial Classification:} Use \textbf{Cross Entropy with Softmax} to ensure probabilities sum to one.
    \item \textbf{Regression:} Generally relies on \textbf{Mean Squared Error (MSE)}.
\end{itemize}

The loss function optimization involves minimizing the loss function by adjusting the weight matrices via a process called \textbf{Gradient Descent}. Note that each layer has its own set of weights, and the loss is a function of the entire network's weights.

\textbf{Gradient Descent Algorithm}
\begin{enumerate}
    \item Initialize weights
    \item Loop until convergence:
    \item Compute gradient, $\frac{\partial J(W)}{\partial W}$
    \item Update weights, $W^{t+1} \leftarrow W^{t} - \eta \frac{\partial J(W)}{\partial W^\top}$
    \item Return weights
\end{enumerate}

We also looked at the following questions and there answers:
\begin{enumerate}
    \item \textbf{How to set initial weights?} \\
    Weights are initialized randomly to break symmetry and ensure the network does not get stuck in a local minimum. There are many different ways to randomly initialize weights, but the most common method uses a Gaussian distribution with a mean of zero.

    \item \textbf{How to compute gradient?} \\
    To determine the optimal weights, we must calculate the gradient. This process involves a \textbf{forward pass} to compute the loss, followed by \textbf{backpropagation} to calculate gradients via the chain rule. This is repeated for every weight in the network using gradients derived from later layers.

    \item \textbf{How to choose the stepsize?} \\
    Weights are updated in the opposite direction of the gradient by a specific step size (learning rate). An appropriate step size helps avoid getting stuck in a local minimum (since the loss function is non-convex and has multiple minima).
    
    \textit{Note on Optimization Speed:} Calculating gradients over an entire dataset is computationally expensive. To speed this up, we often use \textbf{Stochastic Gradient Descent (SGD)} or \textbf{Mini-Batch Gradient Descent}. Additionally, we use adaptive learning rate algorithms such as \textbf{Adagrad}, \textbf{RMSprop}, or \textbf{Adam} to manage learning rates effectively.

    \item \textbf{When to stop?} \\
    We typically run the algorithm until convergence. However, over-optimization can cause the model to overfit the training data. To prevent this, we employ regularization techniques such as \textbf{Norm-Regularization}, \textbf{Dropout}, and \textbf{Early Stopping} to discourage overly complex models and improve generalization on unseen data.
\end{enumerate}

\subsection*{Honor Code}
I pledge my honor that this lecture summary is my own work and adheres to the guidelines in the instructions.
- Daniyal Shahid Iqbal

\end{document}